% Secao 1 - Motivacao e o gancho. Abre com a tensao economica, nao com o SUS.

\section{Motivação: dois mapas para o mesmo país}
\label{sec:motivacao}

Políticas de acesso à saúde são desenhadas sobre mapas. Programas de combate a
desertos de maternidade, a regionalização do SUS em regiões de saúde, a alocação
de novos leitos --- todos perguntam, no fundo, a mesma coisa: quão longe está o
hospital mais próximo? A distância é a régua natural porque é observável,
barata e intuitiva. O problema é que o paciente não viaja em linha reta ao ponto
mais próximo. Ele viaja para onde há o cuidado de que precisa, e esse destino
revelado frequentemente contradiz o mapa.

Este relatório leva essa contradição a sério. Ele contrasta dois conceitos de
deserto hospitalar para os \valNmunic{} municípios brasileiros entre 2015 e 2022. O
primeiro é o deserto de \emph{distância}: a quilometragem ao hospital ativo mais
próximo, a medida que a política baseada em mapa enxerga. O segundo é o deserto
de \emph{fluxo}: o burden efetivo de viagem, ponderado pelos hospitais em que os
residentes de fato internam. Os dois compartilham o suficiente para medir o
mesmo fenômeno e divergem o suficiente para que a escolha da régua mude quem é
classificado como desassistido.

A pergunta é deliberadamente descritiva, não causal: \emph{onde, e para quem, os
dois conceitos discordam?} A divergência não é pequena nem aleatória. Existe um
contingente grande de municípios que o mapa de distância declara atendidos ---
têm um hospital por perto --- mas cujos pacientes, na prática, viajam longe
porque o hospital local não resolve o que precisam. Há também o inverso, na
Amazônia, onde a distância rodoviária exagera um isolamento que o fluxo real, por
rio e local, desmente. Medir acesso por quilômetros, mostramos, não é uma
aproximação inócua: é um mapa diferente, que erra de forma sistemática e
identificável. A questão de se essa diferença custa vidas é real, mas exige
identificação causal e fica explicitamente fora do escopo deste relatório
(Seção~\ref{sec:mortalidade}); aqui, estabelecemos a medida e a geografia do
desacordo.
